\section*{Phụ lục}
\addcontentsline{toc}{section}{Phụ lục}

\subsection*{Phụ lục A: Cấu trúc cơ sở dữ liệu khởi tạo (SQL Schema Migration)}
Dưới đây là mã lệnh SQL (PostgreSQL 16 với phần mở rộng \texttt{pgvector}) dùng để khởi tạo cấu trúc bảng và chỉ mục cho nền tảng SaaS theo kiến trúc Vi dịch vụ:

\begin{lstlisting}[language=SQL, caption={Mã lệnh khởi tạo DDL (Data Definition Language) cho hệ thống}]
-- 1. Kich hoat phan mo rong pgvector cho truy xuat RAG
CREATE EXTENSION IF NOT EXISTS vector;
CREATE EXTENSION IF NOT EXISTS "uuid-ossp";

-- 2. Bang bang thong tin va trang thai cua hang F&B (Tenants)
CREATE TABLE IF NOT EXISTS tenants (
    tenant_id SERIAL PRIMARY KEY,
    name VARCHAR(255) NOT NULL,
    status VARCHAR(50) DEFAULT 'ACTIVE' CHECK (status IN ('ACTIVE', 'LOCKED', 'MAINTENANCE')),
    total_orders BIGINT DEFAULT 0,
    hotline VARCHAR(50),
    telegram_chat_id VARCHAR(100),
    created_at TIMESTAMP WITH TIME ZONE DEFAULT CURRENT_TIMESTAMP,
    updated_at TIMESTAMP WITH TIME ZONE DEFAULT CURRENT_TIMESTAMP
);

-- 3. Bang lich su giao dich don hang (Order Logs)
CREATE TABLE IF NOT EXISTS order_logs (
    order_id UUID PRIMARY KEY DEFAULT uuid_generate_v4(),
    tenant_id INTEGER NOT NULL REFERENCES tenants(tenant_id) ON DELETE CASCADE,
    order_code VARCHAR(100) NOT NULL,
    customer_name VARCHAR(255),
    customer_phone VARCHAR(50),
    total_amount BIGINT NOT NULL CHECK (total_amount >= 0),
    status VARCHAR(50) DEFAULT 'COMPLETED',
    created_at TIMESTAMP WITH TIME ZONE DEFAULT CURRENT_TIMESTAMP
);

-- 4. Bang thuc don va vector ngu nghia (Menu Items & Embeddings)
CREATE TABLE IF NOT EXISTS menu_items (
    item_id SERIAL PRIMARY KEY,
    tenant_id INTEGER NOT NULL REFERENCES tenants(tenant_id) ON DELETE CASCADE,
    item_name VARCHAR(255) NOT NULL,
    price BIGINT NOT NULL,
    description TEXT,
    embedding VECTOR(384) NOT NULL,
    is_available BOOLEAN DEFAULT TRUE,
    created_at TIMESTAMP WITH TIME ZONE DEFAULT CURRENT_TIMESTAMP
);

-- 5. Tao cac chi muc (Indexes) toi uu hieu nang truy van
CREATE INDEX IF NOT EXISTS idx_tenants_status ON tenants(tenant_id, status);
CREATE INDEX IF NOT EXISTS idx_orders_tenant ON order_logs(tenant_id, created_at DESC);
CREATE INDEX IF NOT EXISTS idx_menu_vectors ON menu_items 
USING ivfflat (embedding vector_cosine_ops) WITH (lists = 100);

-- 6. Du lieu mau (Seed Data) kiem thu he thong
INSERT INTO tenants (tenant_id, name, status, hotline) 
VALUES 
(101, 'Tiệm Cơm Sườn Bì Chả Sài Gòn', 'ACTIVE', '0901234567'),
(102, 'Chuỗi Cà Phê AI RAG Express', 'ACTIVE', '0909888999')
ON CONFLICT (tenant_id) DO NOTHING;
\end{lstlisting}

\subsection*{Phụ lục B: Kịch bản kiểm thử tải trọng Locust (\texttt{locustfile.py})}
Mã nguồn chương trình kiểm thử hiệu năng bằng Locust, giả lập 500 khách hàng liên tục gửi đơn đặt hàng và tra cứu thực đơn với lưu lượng đạt đỉnh $> 1.000$ đơn/phút:

\begin{lstlisting}[language=Python, caption={Kịch bản Python giả lập tải cao gửi về \texttt{omni-gateway-service}}]
import json
import time
import hmac
import hashlib
import random
from locust import HttpUser, task, between, events

SECRET_KEY = b"enterprise_saas_secret_key_2026"

class OmnichannelShopperUser(HttpUser):
    wait_time = between(0.1, 0.5) # Gia lap user nhan tin lien tuc
    
    def _generate_hmac_header(self, payload_dict: dict) -> dict:
        body_bytes = json.dumps(payload_dict).encode("utf-8")
        sig = "sha256=" + hmac.new(SECRET_KEY, body_bytes, hashlib.sha256).hexdigest()
        return {"Content-Type": "application/json", "X-Hub-Signature": sig}
        
    @task(3)
    def simulate_menu_inquiry(self):
        """Gia lap luong hoi dap thuc don AI RAG"""
        questions = [
            "Quán có món cơm sườn nào không cay cho trẻ em không?",
            "Cho mình xin bảng giá đồ uống lạnh buổi sáng",
            "Món bán chạy nhất của quán hôm nay là gì?"
        ]
        payload = {
            "tenant_id": random.choice([101, 102]),
            "sender_id": f"fb_user_{random.randint(1000, 9999)}",
            "message_type": "INQUIRY",
            "content": random.choice(questions)
        }
        headers = self._generate_hmac_header(payload)
        
        with self.client.post("/api/v1/webhook/ingest", json=payload, headers=headers, catch_response=True) as resp:
            if resp.status_code in [200, 202]:
                resp.success()
            else:
                resp.failure(f"Inquiry failed with status {resp.status_code}")

    @task(7)
    def simulate_order_submission(self):
        """Gia lap luong chot don hang vao RabbitMQ (> 1000 don/phut)"""
        order_payload = {
            "order_code": f"ORD-{int(time.time())}-{random.randint(10, 99)}",
            "items": [{"name": "Cơm sườn đặc biệt", "qty": 2, "price": 50000}],
            "total_amount": 100000,
            "customer_phone": "0901234567"
        }
        payload = {
            "tenant_id": 101,
            "sender_id": f"zalo_user_{random.randint(10000, 99999)}",
            "message_type": "ORDER",
            "content": json.dumps(order_payload)
        }
        headers = self._generate_hmac_header(payload)
        
        with self.client.post("/api/v1/webhook/ingest", json=payload, headers=headers, catch_response=True) as resp:
            if resp.status_code == 202:
                resp.success()
            else:
                resp.failure(f"Order submission failed: {resp.status_code}")
\end{lstlisting}